\chapter{Глобальная масштабная ковариантность и absolute-scale no-go}

До построения radion-потенциала необходимо отделить физическое уравнение
масштаба от выбора единиц. Пусть вся метрика компактного носителя
масштабируется гомотетией
\begin{equation}
 g\longmapsto \lambda^2 g,
 \qquad \lambda>0.
\end{equation}
Топология, фундаментальная группа, спин-структуры и конечная алгебра при
этом не меняются.

\section{Масштабирование геометрических данных}

Для \(d\)-мерного фактора выполняется
\begin{align}
 L&\longmapsto\lambda L,\\
 \Vol_d&\longmapsto\lambda^d\Vol_d,\\
 \Delta&\longmapsto\lambda^{-2}\Delta,\\
 D&\longmapsto\lambda^{-1}D.
\end{align}
Для \(p\)-формы с неизменными координатными коэффициентами
\begin{equation}
 \|\omega_p\|_{L^2}^2
 \longmapsto
 \lambda^{d-2p}\|\omega_p\|_{L^2}^2.
\end{equation}
Следовательно, все спектральные массы и KK-шаги масштабируются как
\(\lambda^{-1}\).

Если конечный оператор связан с радиусом условием
\begin{equation}
 M=\frac1R,
\end{equation}
то произведение
\begin{equation}
 MR=1
\end{equation}
инвариантно. Это физическое безразмерное утверждение. Отдельные значения
\(M\) и \(R\) без внешней единицы масштаба не являются предсказаниями.

\begin{theorem}[Scale-covariance theorem]
Топологические данные \(K\), конечная алгебра, first-order constraints и
нормированный finite-Dirac потенциал определяют только гомотетический класс
и безразмерные комбинации. Они не могут выбрать абсолютную массу, поскольку
каждое решение порождает семейство
\begin{equation}
 (g,D_F)\longmapsto
 (\lambda^2g,\lambda^{-1}D_F)
\end{equation}
с теми же безразмерными инвариантами.
\end{theorem}

\section{Spectral cutoff}

Стандартный спектральный функционал имеет вид
\begin{equation}
 S_{\mathrm{spec}}=\Tr f(D/\Lambda).
\end{equation}
При одновременном преобразовании
\begin{equation}
 D\longmapsto\lambda^{-1}D,
 \qquad
 \Lambda\longmapsto\lambda^{-1}\Lambda
\end{equation}
отношение \(D/\Lambda\) не меняется. Поэтому \(\Lambda\) является
dimensionful anchor, если не заменена динамическим dilaton/radion полем.
Сам спектральный след не объясняет абсолютное значение \(\Lambda\).

\section{Почему простая radion-подстановка недостаточна}

Одночленный однородный потенциал
\begin{equation}
 V(R)=cR^q
\end{equation}
не имеет изолированного положительного минимума. Минимальный двухчленный
пример
\begin{equation}
 V(R)=AR+\frac{B}{R},
 \qquad A,B>0,
\end{equation}
даёт
\begin{equation}
 R_*=\sqrt{\frac BA}.
\end{equation}
Следовательно, физический радиус определяется отношением двух
коэффициентов. Равенство \(A=B\), приводящее к \(R_*=1\), должно следовать
из независимо объявленной duality или общей квадратной структуры; оно не
следует из топологии окружности.

Аналогично, запись \(R=1\), использованная для объёмов
\(\pi^2\) и \(2\pi\), фиксирует безразмерную репрезентацию геометрии, но не
переводит единицу в ГэВ\(^{-1}\).

\begin{nogocriterion}[Absolute-scale no-go]
В классе parent actions, содержащих только нормированную геометрию \(K\),
конечную алгебру и безразмерный spectral flattening, абсолютная массовая
шкала не выводится. Для её появления необходим хотя бы один из механизмов:
\begin{enumerate}
 \item dimensionful gravitational input;
 \item динамический dilaton/radion с изолированным vacuum;
 \item квантовая dimensional transmutation;
 \item заранее объявленная физическая cutoff-шкала.
\end{enumerate}
\end{nogocriterion}

\section{Следствие для статуса предсказаний}

До закрытия scale gate допустимы только:
\begin{itemize}
 \item безразмерные отношения масс;
 \item углы и фазы;
 \item ранги, индексы и кратности;
 \item отношения связей и normalization-sensitive безразмерные нормы.
\end{itemize}
Абсолютные значения масс могут использоваться лишь как один scale-setting
input, но не как независимые blind-предсказания.

\begin{protocol}[Развилка следующего этапа]
Версия III должна выбрать один из двух честных маршрутов:
\begin{enumerate}
 \item построить совместный dilaton--radion--finite-Dirac action и вывести
 физический scale vacuum;
 \item заморозить один dimensionful input и проводить blind-тест только на
 безразмерных наблюдаемых, не объявляя абсолютные массы предсказанными.
\end{enumerate}
\end{protocol}

Машинный аудит сохранён в
\path{s2t_v3_absolute_scale_no_go_results.json}.